% Spanish for Beginners - Week 14 Cornell Final Notes (Spring 2026, Wed/Thu)
% Compile: pdflatex spanish-for-beginners-final-cornell-master-notes-wed-thu-spring2026.tex
\documentclass[10pt,a4paper]{article}
\usepackage[utf8]{inputenc}
\usepackage[T1]{fontenc}
\usepackage[margin=1.2cm]{geometry}
\usepackage{xcolor}
\usepackage{longtable}
\usepackage{enumitem}
\setlist[itemize]{leftmargin=1.2em,itemsep=1pt,topsep=2pt}
\definecolor{spblue}{RGB}{20,70,130}
\definecolor{soft}{RGB}{244,248,252}
\setlength{\parindent}{0pt}
\setlength{\parskip}{0.35em}

\newcommand{\cornell}[3]{
\subsection*{#1}
\begin{longtable}{|p{0.22\linewidth}|p{0.74\linewidth}|}
\hline
\rowcolor{soft}\textbf{Cue / Questions} & \textbf{Notes} \\
\hline
#2 & #3 \\
\hline
\end{longtable}
}

\begin{document}
\begin{center}
{\Large\bfseries\color{spblue} Spanish for Beginners - Week 14 Final Cornell Master Notes}\\
{\small Spring 2026 (typical class days: Wed/Thu) | Consolidated from class notes, unit practice, and midterm materials}
\end{center}

\section*{Performance and Focus Rules}
\begin{itemize}
  \item \textbf{Spanish marker:} A- (7.1/10) = notable/bueno (NT).
  \item \textbf{Goal:} preserve high performance while improving consistency in grammar accuracy.
  \item \textbf{Global fail thresholds:} WIP 3.49--0; UCM 4.99--0; suspenso = 4.
\end{itemize}

\section*{Cornell Notes by Core Language Block}

\cornell{1) Present Tense and Core Verb Control}
{How are regular verbs formed?\\
Where do common irregulars appear?\\
What errors repeat most?}
{
\textbf{Regular pattern control:} -ar, -er, -ir endings by person must be automatic.\\
\textbf{Irregular anchors:} ser, estar, tener, ir, hacer, preferir need high-frequency repetition.\\
\textbf{Accuracy rule:} subject-verb agreement is scored quickly in oral/written output.\\
\textbf{Practice strategy:} short daily conjugation loops + immediate sentence application.\\
\textbf{Exam tip:} prioritize grammatical correctness over overly complex sentence length.
}

\cornell{2) Essential Grammar Contrasts}
{Ser vs estar?\\
Por vs para?\\
Direct vs indirect object?}
{
\textbf{Ser/estar logic:} identity/essential trait vs temporary condition/location.\\
\textbf{Por/para logic:} cause/means/exchange vs purpose/recipient/deadline.\\
\textbf{Object pronouns:} correct placement and agreement improve fluency perception.\\
\textbf{Mandates and infinitives:} use clear command structures in practical tasks.\\
\textbf{High-value move:} give one correct contrast example in exam responses.
}

\cornell{3) Unit 4 Functional Communication}
{How to express plans?\\
How to describe clothing and preferences?\\
How to sustain short conversation?}
{
\textbf{Plans and intention:} ir + a + infinitive for near-future communication.\\
\textbf{Preference language:} preferir, gustar-type structures, and comparative phrasing.\\
\textbf{Description domain:} ropa, colores, and possession expressions (es/de).\\
\textbf{Conversation performance:} prioritize clear, correct, and simple turns.\\
\textbf{Scoring insight:} communicative success plus grammatical reliability drives strong marks.
}

\cornell{4) Listening and Oral Response Technique}
{How to improve listening extraction?\\
How to respond under time pressure?\\
What is exam-safe structure?}
{
\textbf{Listening method:} capture keywords (who, where, when, action) before detail.\\
\textbf{Oral structure:} opening phrase -> key info -> short expansion -> close.\\
\textbf{Error control:} avoid long uncertain constructions under pressure.\\
\textbf{Confidence tool:} prepare reusable sentence frames for common prompts.\\
\textbf{Exam readiness:} repetition with timed mini-prompts improves response speed.
}

\cornell{5) Writing and Final Review Routine}
{How to revise efficiently?\\
How to reduce avoidable mistakes?\\
How to keep A- or improve?}
{
\textbf{Writing checklist:} agreement, verb tense consistency, punctuation, accents (when required by course).\\
\textbf{Weekly loop:} vocab review -> grammar drill -> speaking practice -> short writing output.\\
\textbf{Retention mechanism:} active recall and production, not passive rereading only.\\
\textbf{Finals rule:} maintain consistency and avoid ``small grammar losses'' that lower grade band.\\
\textbf{Goal state:} stable NT performance trending toward SB threshold.
}

\noindent\textbf{A1 benchmark (CEFR ``Breakthrough'' / Beginner):} Expectation is basic survival Spanish---introduce yourself, ask directions, describe your immediate surroundings with simple sentences. Blocks 6--14 summarize the usual A1 pillars (grammar, vocab themes, can-do skills, structure, pronunciation, and street-survival action phrases).

\cornell{6) A1 --- Essential Grammar (Mostly Present Tense)}
{Subject pronouns?\\
Ser vs estar?\\
Regular and irregular verbs?\\
Gender and number?}
{
\textbf{Present tense first:} A1 centers on the present (not only, but heavily).\\
\textbf{Pronouns:} \textit{tú} (informal) vs \textit{usted} (formal); \textit{nosotros} vs \textit{nosotras}.\\
\textbf{Two ``to be'' verbs:} \textit{ser} for stable traits (nationality, job, defining description); \textit{estar} for temporary states (mood, health) and location.\\
\textbf{Regular verbs:} -ar, -er, -ir present conjugation must be automatic.\\
\textbf{Key irregulars:} \textit{tener}, \textit{ir}, \textit{hacer} at high frequency.\\
\textbf{Agreement:} noun gender/number; adjectives match (e.g.\ \textit{el gato negro}, \textit{las gatas negras}).
}

\cornell{7) A1 --- Core Vocabulary Themes}
{Which domains build basic sentences?\\
What ``bricks'' repeat in exams and oral tests?}
{
\textbf{Greetings and introductions:} \textit{Hola}, \textit{Buenos días}, \textit{¿Cómo te llamas?}\\
\textbf{Numbers 1--100:} age, phone, prices, quantities.\\
\textbf{Family:} \textit{madre}, \textit{padre}, \textit{hermano}, \textit{hijo}, etc.\\
\textbf{Time and calendar:} clock time, days of the week, months.\\
\textbf{Food and drink:} supermarket/restaurant survival set.\\
\textbf{Home:} rooms and basic furniture (\textit{mesa}, \textit{silla}, \textit{cama}).\\
\textbf{Colors and physical description:} people and objects.
}

\cornell{8) A1 --- Functional Skills (CEFR Can-Do)}
{What must you actually perform?\\
What counts as ``simple interaction''?}
{
\textbf{Introduce self and others:} name, where you live, basic occupation/student status.\\
\textbf{Personal Q\&A:} origin, simple hobbies/preferences.\\
\textbf{Simple interaction:} communication works when the partner speaks slowly and helps---your job is clear, short turns.\\
\textbf{Basic shopping:} ask price (\textit{¿Cuánto cuesta?}); recognize numbers/prices in reply.
}

\cornell{9) A1 --- Sentence Structure (vs English Habits)}
{Adjective placement?\\
Negation?\\
Questions?}
{
\textbf{SVO familiarity:} basic order often parallels English; focus on exceptions.\\
\textbf{Adjectives usually after the noun:} \textit{La casa azul.} (not ``azul casa'').\\
\textbf{Negation:} \textit{no} + verb (\textit{No hablo japonés.}).\\
\textbf{Questions:} opening/closing marks \textit{¿\ldots?} (and \textit{¡\ldots!} for emphasis); learn interrogatives---\textit{qué}, \textit{quién}, \textit{dónde}, \textit{por qué}, \textit{cómo}.
}

\cornell{10) A1 --- Pronunciation Basics}
{Which letters trip English speakers?\\
Why is spelling ``phonetic''?}
{
\textbf{Phonetic default:} letters usually match sounds consistently (few silent traps beyond \textit{h}).\\
\textbf{H:} always silent (\textit{hola} $\approx$ ``ola'').\\
\textbf{J:} breathy/velar friction (not English ``j'' as in ``jam'').\\
\textbf{LL:} often like English ``y'' in many lects (\textit{pollo}).\\
\textbf{Ñ:} palatal nasal (\textit{España}) like ``ny'' in ``canyon''.
}

\cornell{11) A1 Street Survival --- Action Verbs}
{Present-tense “use it now” verbs?\\
What do you say when you want/need something?}
{
\textbf{Quiero:} “I want” (e.g.\ \textit{Quiero un café.}).\\
\textbf{Necesito:} “I need” (e.g.\ \textit{Necesito el baño.}).\\
\textbf{Tengo:} “I have / I’m” (age: \textit{Tengo 20 años}).\\
\textbf{Compro:} “I buy” (survival shopping verb).\\
\textbf{Entiendo:} “I understand” (useful when responding slowly/clearly).
}

\cornell{12) A1 Street Survival --- Directions}
{How do you ask for places?\\
How do you guide someone?}
{
\textbf{¿Dónde está \ldots?}:} “Where is \ldots?”\\
\textbf{A la derecha / a la izquierda:} to the right / to the left.\\
\textbf{Todo recto:} straight ahead.\\
\textbf{Cerca / Lejos:} near / far.\\
\textbf{La entrada / La salida:} the entrance / the exit.
}

\cornell{13) A1 Street Survival --- Shopping & Eating Out}
{How do you order and pay?\\
How do you handle simple restrictions?}
{
\textbf{La cuenta, por favor:} the bill, please.\\
\textbf{¿Cuánto cuesta?:} how much does it cost?\\
\textbf{Una mesa para dos:} a table for two.\\
\textbf{Soy alérgico/a a \ldots:} I am allergic to \ldots\\
\textbf{Sin \ldots (cebolla/carne):} without \ldots (onion/meat).
}

\cornell{14) A1 Street Survival --- High-frequency Adjectives}
{How do you describe size, cost, quality, and temperature?\\
Where do adjectives go?}
{
\textbf{Grande / pequeño:} big / small.\\
\textbf{Caro / barato:} expensive / cheap.\\
\textbf{Bueno / malo:} good / bad.\\
\textbf{Caliente / frío:} hot / cold.\\
\textbf{Bonito / feo:} pretty / ugly.\\
\textbf{Grammar move:} adjectives go AFTER the noun (\textit{un café barato}).
}

\section*{Finals Tile Insights (From Uploaded Images)}
\begin{itemize}
  \item The exam tiles show strongest gains come from \textbf{grammar precision}: interrogatives, agreement, and preposition usage.
  \item Listening/reading tasks improve when answers stay literal and avoid over-guessing details.
  \item Written prompt feedback suggests keeping sentence structure simple but correct under time pressure.
  \item Unit 4 vocabulary and clothing/travel dialogue prompts are clearly relevant for final revision.
\end{itemize}

\section*{Summary Block (Cornell Bottom)}
\begin{itemize}
  \item Map course revision onto A1 ``can-do'' survival goals (blocks 6--14) while drilling blocks 1--5 course priorities.
  \item Keep Spanish in high band by protecting grammar precision and response clarity.
  \item Use concise, correct sentences during finals; complexity only when confidence is high.
  \item Maintain daily micro-practice to preserve speed and accuracy.
\end{itemize}

\end{document}
